\chapter{Spectral Graph Theory}

A graph is a combinatorial object: a set of vertices, a set of edges, no numbers anywhere unless you put them there. Yet the moment you write the graph down as a matrix -- the adjacency matrix, or better, the Laplacian -- you have handed it to linear algebra, and linear algebra hands back a list of real numbers, the \emph{eigenvalues}. Those numbers turn out to encode an astonishing amount of structure: whether the graph is connected, how many pieces it splits into, whether it is bipartite, how well-knit it is, and -- the deepest fact in this chapter -- whether it can be cut into two large pieces by severing only a few edges. Spectral graph theory is the study of the dictionary between the combinatorics of $\Gph$ and the spectrum of its matrices.

The previous chapter (Chapter~31) built the Laplacian $\Lap = \mathbf{D} - \Adjm$ and proved its basic identity, the quadratic form $x^\top \Lap x = \sum_{(i,j)\in E}(x_i - x_j)^2$. We lean on that identity constantly here. The payoff of the present chapter is twofold. First, we read graph structure straight off the spectrum: the multiplicity of the eigenvalue $0$ counts connected components, the largest eigenvalue of the normalized Laplacian touches $2$ exactly when a component is bipartite, and the \emph{second} eigenvalue -- the algebraic connectivity, or Fiedler value -- measures how robustly the graph hangs together. Second, we prove the \emph{Cheeger inequality}, the theorem that ties that second eigenvalue to the existence of a sparse cut, and whose constructive direction is the engine behind spectral clustering (Chapter~33) and the analysis of expanders (Chapter~34).

We keep everything self-contained and we trace a real little graph by hand. By the end you should be able to take a six-vertex graph, write its Laplacian, find its spectrum, point at the Fiedler vector, and read off the near-optimal cut it suggests.

\begin{historybox}
Spectral graph theory bridges continuous differential geometry and discrete computer science.
In 1970, American mathematician Jeff Cheeger was investigating Riemannian manifolds, asking:
how does the lowest non-zero eigenvalue $\lambda_1$ of the continuous Laplace--Beltrami operator
relate to the narrowest bottleneck or geometric constriction of the manifold? Cheeger proved
the fundamental continuous inequality $h(M) \ge \tfrac{\lambda_1}{2}$, tying eigenvalues
directly to isoperimetric geometry.

Simultaneously, computer scientists designing VLSI circuits and telephone switching networks
faced the Sparsest Cut and Minimum Bisection problems: dividing a circuit graph into two
balanced halves while severing as few interconnect wires as possible. Sparsest cut is NP-hard,
and early heuristic searches (such as the Kernighan--Lin algorithm) constantly became trapped
in poor local minima. In 1972, William Donath and Alan Hoffman proposed relaxing the discrete
binary cut assignment $x_i \in \{-1, +1\}$ to a continuous real vector $x \in \R^n$.
Miroslav Fiedler (1973) recognized that this continuous quadratic relaxation is minimized
precisely by the second Laplacian eigenvector $v_2$ (the Fiedler vector)!

Yet a fundamental controversy lingered: because $v_2$ is a continuous real vector, could its
coordinates be misleading? What guarantee existed that rounding $v_2$ to a discrete cut
would yield a cut anywhere near the true combinatorial optimum?
In 1985, Noga Alon and Vitali Milman (and independently Alistair Sinclair and Mark Jerrum in
1989 in the context of Markov chain mixing times) proved the \emph{discrete Cheeger inequality}:
\[
  \tfrac{\mu_2}{2} \;\le\; \phi(\Gph) \;\le\; \sqrt{2\,\mu_2} .
\]
They proved that sorting vertices by their Fiedler vector coordinates and evaluating the $n-1$
sweep cuts provably guarantees a cut satisfying this quadratic bound. This transformed spectral
graph theory from a heuristic intuition into a mathematically rigorous foundation for graph
partitioning and algorithm design.
\end{historybox}

\chapterline

\section{Setup: which matrix, which spectrum}
\label{ch32:setup}

Fix an undirected graph $\Gph=(V,E)$ with $n=\abs{V}$ vertices and $m=\abs{E}$ edges, possibly weighted with nonnegative weights $w_{ij}\ge 0$ on edge $(i,j)$. Let $d_i=\sum_j w_{ij}$ be the (weighted) degree of vertex $i$ and $\mathbf{D}=\operatorname{diag}(d_1,\dots,d_n)$. We will use three matrices.

\begin{definition}[The three matrices]
\label{ch32:three-matrices}
The \emph{adjacency matrix} $\Adjm$ has $\Adjm_{ij}=w_{ij}$ (and $0$ on the diagonal, no self-loops). The \emph{combinatorial Laplacian} is
\[
  \Lap \;=\; \mathbf{D} - \Adjm .
\]
The \emph{normalized Laplacian} (assuming no isolated vertices, so every $d_i>0$) is
\[
  \mathcal{L} \;=\; \mathbf{D}^{-1/2}\,\Lap\,\mathbf{D}^{-1/2} \;=\; \mathbf{I} - \mathbf{D}^{-1/2}\Adjm\,\mathbf{D}^{-1/2}.
\]
\end{definition}

Both $\Lap$ and $\mathcal{L}$ are real symmetric, so by the spectral theorem each has $n$ real eigenvalues and an orthonormal basis of real eigenvectors. We order the Laplacian eigenvalues
\[
  0 = \lambda_1 \le \lambda_2 \le \cdots \le \lambda_n
\]
and the normalized ones
\[
  0 = \mu_1 \le \mu_2 \le \cdots \le \mu_n \le 2 .
\]
(We will earn each of these endpoints below; that $\lambda_1=0$ and $\mu_1=0$ comes from the all-ones vector, and $\mu_n\le 2$ is Lemma~\ref{ch32:mu-le-2}.)

\begin{intuitionbox}
Think of a vector $x\in\R^V$ as an assignment of a real number to each vertex -- a ``signal'' on the graph. The Laplacian quadratic form
\[
  x^\top \Lap x \;=\; \sum_{(i,j)\in E} w_{ij}\,(x_i-x_j)^2
\]
measures how \emph{rough} the signal is: it is large when the signal jumps across edges and zero when the signal is constant on each connected piece. Eigenvectors of small eigenvalue are the smoothest non-constant signals the graph admits. ``Smooth'' here means: neighbors get nearly equal values. The whole chapter is the consequences of that one sentence.
\end{intuitionbox}

Why two Laplacians? The combinatorial $\Lap$ is the natural one for many proofs and for the matrix-tree theorem; the normalized $\mathcal{L}$ is the right one for cuts and random walks, because it weights each vertex by its degree. When the graph is $d$-regular (every $d_i=d$) the two coincide up to scaling: $\mathcal{L}=\frac1d\Lap$, so $\mu_k=\lambda_k/d$. We will state the structural facts for whichever matrix makes them cleanest and translate when it matters.

\begin{focusbox}
Three questions organize the chapter.
\begin{enumerate}
  \item \textbf{Structure from the spectrum.} What do $\lambda_1,\lambda_n$, multiplicities, and the bipartite case tell us about $\Gph$? (Section~\ref{ch32:structure})
  \item \textbf{The second eigenvalue.} Why is $\lambda_2$ -- the algebraic connectivity, the Fiedler value -- the single most informative number, and what does its eigenvector look like? (Section~\ref{ch32:fiedler})
  \item \textbf{Cheeger.} How does $\lambda_2$ (or $\mu_2$) sandwich the sparsest cut, and how do we turn the eigenvector into an actual cut? (Section~\ref{ch32:cheeger})
\end{enumerate}
\end{focusbox}

\subsection{The variational principles we will reuse}

Everything spectral here flows from two standard facts about symmetric matrices. We state them once and cite them by name.

\begin{lemma}[Rayleigh quotient and the bottom of the spectrum]
\label{ch32:rayleigh}
Let $M$ be symmetric with eigenvalues $\nu_1\le\cdots\le\nu_n$ and orthonormal eigenvectors $v_1,\dots,v_n$. The \emph{Rayleigh quotient} is $R_M(x)=\dfrac{x^\top M x}{x^\top x}$ for $x\ne 0$. Then
\[
  \nu_1=\min_{x\ne 0} R_M(x),\qquad
  \nu_k=\min_{\substack{x\ne 0\\ x\perp v_1,\dots,v_{k-1}}} R_M(x),
\]
and the minimizer of the $k$-th problem is $v_k$.
\end{lemma}

\begin{lemma}[Courant--Fischer min--max]
\label{ch32:courant-fischer}
With the same notation,
\[
  \nu_k \;=\; \min_{\dim S = k}\ \max_{x\in S\setminus\{0\}} R_M(x)
        \;=\; \max_{\dim S = n-k+1}\ \min_{x\in S\setminus\{0\}} R_M(x),
\]
the optimum over $k$-dimensional subspaces $S\subseteq\R^n$.
\end{lemma}

Both are proved in any linear algebra text and in the companion Linear Algebra book; we use them as tools. The first gives us $\lambda_2$ as a constrained minimum of $\sum (x_i-x_j)^2$ over signals orthogonal to the constant vector. The second lets us compare eigenvalues across edge insertions and across the two Laplacians without ever computing them.

\begin{connectionbox}
The Rayleigh-quotient view is exactly the link to the Linear Algebra book's treatment of symmetric matrices and the min--max characterization, and to Chapter~31's identity. It is also the bridge to Chapter~35 (random walks and electrical networks): the eigenvectors of $\mathcal{L}$ are, up to the $\mathbf{D}^{1/2}$ change of variables, the eigenvectors of the random-walk transition matrix $\mathbf{P}=\mathbf{D}^{-1}\Adjm$, and the quadratic form is the energy dissipated by a unit current. Keep that triangle -- Laplacian, random walk, electrical energy -- in view.
\end{connectionbox}

\section{Reading graph structure off the spectrum}
\label{ch32:structure}

\subsection{The kernel: connectivity and components}

\begin{theorem}[Multiplicity of $0$ counts components]
\label{ch32:components}
For the combinatorial Laplacian $\Lap$ of an undirected graph, the multiplicity of the eigenvalue $0$ equals the number $c$ of connected components of $\Gph$. The same holds for $\mathcal{L}$ (with no isolated vertices). In particular $\lambda_2>0$ if and only if $\Gph$ is connected.
\end{theorem}

\begin{proof}
$\Lap$ is positive semidefinite because $x^\top\Lap x=\sum_{(i,j)\in E}w_{ij}(x_i-x_j)^2\ge 0$ for all $x$; hence every eigenvalue is $\ge 0$, and $\lambda=0$ has eigenvectors exactly the $x$ with $x^\top\Lap x=0$, i.e. the kernel of $\Lap$.

Now $x^\top\Lap x=0$ forces $x_i=x_j$ for every edge $(i,j)$ with $w_{ij}>0$. Walking along edges, $x$ is constant on each connected component. Conversely any vector that is constant on each component makes every term vanish and so lies in the kernel. The indicator vectors $\mathbf{1}_{C_1},\dots,\mathbf{1}_{C_c}$ of the $c$ components are therefore a basis of the kernel: they are linearly independent (by disjoint coordinate supports) and they span it (any kernel vector is a combination of them). Hence $\dim\ker\Lap=c$, which is the multiplicity of $0$. For $\mathcal{L}$, note $\mathcal{L}=\mathbf{D}^{-1/2}\Lap\mathbf{D}^{-1/2}$ is a congruence by an invertible matrix, so it has the same rank and the same number of zero eigenvalues; the kernel vectors are $\mathbf{D}^{1/2}\mathbf{1}_{C_\ell}$.
\end{proof}

The all-ones vector $\mathbf{1}$ is always in the kernel of $\Lap$ (each row sums to $0$), which is why $\lambda_1=0$ for every graph. The content of the theorem is that $0$ has \emph{no other} eigenvectors precisely when the graph is connected.

\begin{trapbox}
A graph with $\lambda_2$ very small but positive is connected, yet only barely -- it has a near-disconnection, a bottleneck. Beginners read $\lambda_2=0$ as ``disconnected'' and stop there. The interesting regime is $\lambda_2$ tiny but nonzero: the graph is one piece, but a thin one, and the Fiedler vector points right at the weak seam. That is the whole motivation for Cheeger.
\end{trapbox}

\subsection{The top of the spectrum and bipartiteness}

For the normalized Laplacian the top eigenvalue $\mu_n$ has its own clean meaning.

\begin{lemma}[$\mu_n\le 2$, with equality for bipartite]
\label{ch32:mu-le-2}
For the normalized Laplacian, $0\le\mu_k\le 2$ for all $k$. Moreover, on a connected graph, $\mu_n=2$ if and only if $\Gph$ is bipartite.
\end{lemma}

\begin{proof}
Write $y=\mathbf{D}^{-1/2}x$ so that, with the substitution, the normalized Rayleigh quotient becomes
\[
  \frac{x^\top\mathcal{L}x}{x^\top x}
   = \frac{\sum_{(i,j)\in E} w_{ij}(y_i-y_j)^2}{\sum_i d_i\, y_i^2}.
\]
By $(y_i-y_j)^2\le 2(y_i^2+y_j^2)$ and $\sum_{(i,j)} w_{ij}(y_i^2+y_j^2)=\sum_i d_i y_i^2$, the numerator is at most $2$ times the denominator, so $\mu_n\le 2$. Equality needs $(y_i+y_j)^2=0$ on every edge, i.e. $y_i=-y_j$ across each edge with a consistent sign; this two-colors the vertices, so the graph is bipartite. Conversely if $\Gph$ is bipartite with sides $A,B$, take $y_i=+1$ on $A$ and $y_i=-1$ on $B$; then every $(y_i-y_j)^2=4=2(y_i^2+y_j^2)$, achieving the bound, so $\mu_n=2$.
\end{proof}

\begin{intuitionbox}
$\mu_n=2$ is the spectral fingerprint of bipartiteness: there is a signal that flips sign across \emph{every} edge, the roughest possible signal. ``Almost bipartite'' graphs have $\mu_n$ close to $2$ but below it. So the normalized spectrum brackets the graph from both ends: $\mu_2$ small says ``almost disconnected''; $2-\mu_n$ small says ``almost bipartite.'' Both are kinds of weak structure, and both are caught spectrally.
\end{intuitionbox}

\subsection{A first worked spectrum: the path $P_4$}
\label{ch32:p4}

Let us pin all of this to a concrete graph before going further. Take the path on four vertices, $P_4$: edges $1$--$2$, $2$--$3$, $3$--$4$, unit weights.

\begin{center}
\begin{tikzpicture}[every node/.style={circle,draw=nodeblue,thick,fill=blue!6,minimum size=8mm,inner sep=0pt},
  edge/.style={draw=edgegray,very thick}]
  \node (v1) at (0,0) {$1$};
  \node (v2) at (2,0) {$2$};
  \node (v3) at (4,0) {$3$};
  \node (v4) at (6,0) {$4$};
  \draw[edge] (v1)--(v2)--(v3)--(v4);
\end{tikzpicture}
\end{center}

The degrees are $1,2,2,1$, so
\[
  \Lap(P_4) \;=\;
  \begin{pmatrix}
    \phantom{-}1 & -1 & \phantom{-}0 & \phantom{-}0\\
    -1 & \phantom{-}2 & -1 & \phantom{-}0\\
    \phantom{-}0 & -1 & \phantom{-}2 & -1\\
    \phantom{-}0 & \phantom{-}0 & -1 & \phantom{-}1
  \end{pmatrix}.
\]
The eigenvalues of the path Laplacian are known in closed form: $\lambda_k = 2 - 2\cos\!\big(\tfrac{(k-1)\pi}{n}\big)$ for $k=1,\dots,n$. With $n=4$:
\[
  \lambda_1 = 0,\quad
  \lambda_2 = 2-2\cos\tfrac{\pi}{4}=2-\sqrt2\approx 0.586,\quad
  \lambda_3 = 2,\quad
  \lambda_4 = 2+\sqrt2\approx 3.414.
\]
The Fiedler vector (eigenvector for $\lambda_2$) is $v_2=(\cos\tfrac{\pi}{8},\cos\tfrac{3\pi}{8},\cos\tfrac{5\pi}{8},\cos\tfrac{7\pi}{8})\approx(0.92,0.38,-0.38,-0.92)$. Notice it is monotone along the path and changes sign exactly once, between vertices $2$ and $3$ -- which is precisely where you would cut $P_4$ to split it into two equal halves. That sign change is not a coincidence; it is the seed of spectral partitioning, and for trees one can prove the Fiedler vector has exactly one sign change (its zero set is connected).

\begin{labbox}
Confirm $\lambda_2=2-\sqrt2$ by plugging $v_2$ into $\Lap v_2=\lambda_2 v_2$ for, say, the first coordinate: row $1$ gives $v_{2,1}-v_{2,2}=0.92-0.38=0.54$, and $\lambda_2 v_{2,1}=0.586\cdot 0.92\approx 0.54$. Match. Now do the same hand-check for the second coordinate (row $2$): you should get $-v_{2,1}+2v_{2,2}-v_{2,3}=-0.92+0.76+0.38=0.22$ against $0.586\cdot 0.38\approx 0.22$.
\end{labbox}

\subsection{Two more reference spectra}

Memorizing a handful of spectra builds intuition fast.

\textbf{The complete graph $K_n$.} Here $\Lap=n\mathbf{I}-\mathbf{J}$ where $\mathbf{J}$ is the all-ones matrix. So $\Lap\mathbf{1}=0$, and on $\mathbf{1}^\perp$ we have $\mathbf{J}x=0$, hence $\Lap x=nx$. The spectrum is
\[
  \lambda_1=0,\qquad \lambda_2=\lambda_3=\cdots=\lambda_n=n.
\]
The algebraic connectivity is as large as it can be, $\lambda_2=n$: $K_n$ is the most robustly connected graph on $n$ vertices, with no weak cut anywhere.

\textbf{The cycle $C_n$.} By circulant diagonalization the eigenvalues are $\lambda_k=2-2\cos\!\big(\tfrac{2\pi(k-1)}{n}\big)$, $k=1,\dots,n$. For $C_4$: $0,2,2,4$. The algebraic connectivity $\lambda_2=2-2\cos\tfrac{2\pi}{n}\approx \tfrac{4\pi^2}{n^2}$ for large $n$ -- it decays like $1/n^2$, reflecting that a long cycle is easy to cut (snip any two edges).

\begin{historybox}
The systematic study of graph spectra goes back to chemistry and to the theoretical physics of molecular orbitals in the 1930s--50s (the adjacency eigenvalues of a molecule's bond graph approximate $\pi$-electron energies). On the mathematical side, Mark Kac's 1966 question ``Can one hear the shape of a drum?'' is the continuous cousin: do the Laplacian eigenvalues determine the domain? For graphs the answer is famously \emph{no} in general -- there exist \emph{cospectral} non-isomorphic graphs (the smallest pair has $6$ vertices) -- so the spectrum is a powerful but lossy invariant. Miroslav Fiedler's 1973 papers isolated $\lambda_2$ and its eigenvector as the key to connectivity and partitioning, which is why both bear his name.
\end{historybox}

\section{The second eigenvalue: algebraic connectivity}
\label{ch32:fiedler}

\begin{definition}[Algebraic connectivity / Fiedler value]
\label{ch32:def-fiedler}
The \emph{algebraic connectivity} of $\Gph$ is $a(\Gph)=\lambda_2(\Lap)$, the second-smallest Laplacian eigenvalue. A unit eigenvector for $\lambda_2$ is a \emph{Fiedler vector}.
\end{definition}

By Lemma~\ref{ch32:rayleigh} and the fact that $v_1=\mathbf{1}/\sqrt n$,
\begin{equation}
\label{ch32:lambda2-var}
  \lambda_2 \;=\; \min_{\substack{x\ne 0\\ x\perp\mathbf{1}}}\ \frac{\sum_{(i,j)\in E} w_{ij}(x_i-x_j)^2}{\sum_i x_i^2}.
\end{equation}
Read this aloud: among all signals whose entries sum to zero -- the ``balanced'' signals -- $\lambda_2$ is the smallest possible roughness per unit energy. A graph with a thin bottleneck admits a balanced signal that is nearly constant on each side of the bottleneck and only changes across the few cut edges; such a signal has tiny numerator and order-$1$ denominator, forcing $\lambda_2$ small. This is exactly the intuition Cheeger will make quantitative.

\subsection{Why $\lambda_2$ deserves the name ``connectivity''}

\begin{proposition}[Monotonicity and basic bounds]
\label{ch32:monotone}
Algebraic connectivity behaves like a connectivity measure should.
\begin{enumerate}
  \item \emph{(Edge monotonicity.)} Adding an edge (or increasing a weight) never decreases any $\lambda_k$; in particular $\lambda_2$ can only grow. Removing an edge can only shrink it.
  \item \emph{(Bounded by vertex connectivity.)} $\lambda_2(\Gph)\le \kappa(\Gph)$, the vertex connectivity (minimum number of vertices whose deletion disconnects $\Gph$), for $\Gph\ne K_n$. Hence $\lambda_2\le \delta(\Gph)$, the minimum degree.
  \item \emph{(Subgraph by deletion.)} If $\Gph'$ is $\Gph$ with a vertex deleted, $\lambda_2(\Gph')\ge \lambda_2(\Gph)-1$.
\end{enumerate}
\end{proposition}

\begin{proof}[Proof of (1)]
Adding edge $e=(p,q)$ with weight $w$ replaces $\Lap$ by $\Lap + w\,\Lap_e$, where $\Lap_e=(\mathbf{e}_p-\mathbf{e}_q)(\mathbf{e}_p-\mathbf{e}_q)^\top$ is itself positive semidefinite. For any subspace $S$ and any $x\in S$, $x^\top(\Lap+w\Lap_e)x\ge x^\top\Lap x$, so the Courant--Fischer max over $S$ only increases; taking the min over $S$ gives $\lambda_k(\Lap+w\Lap_e)\ge\lambda_k(\Lap)$. (This is Weyl monotonicity for a PSD perturbation.)
\end{proof}

Part (1) is the cleanest evidence that $\lambda_2$ is a connectivity. More edges $\Rightarrow$ harder to disconnect $\Rightarrow$ larger $\lambda_2$, and the inequality is exact and one-line. Part (2), Fiedler's original bound $\lambda_2\le\kappa$, says the spectral measure never overstates how connected the graph is; we omit its proof (a clean argument tests the Rayleigh quotient on a vector supported away from a vertex cut).

\subsection{The Fiedler vector as an embedding}

Equation~\eqref{ch32:lambda2-var} says the Fiedler vector $v_2$ places each vertex on the real line so that adjacent vertices stay close (small $\sum(x_i-x_j)^2$) while the configuration stays spread out and balanced ($x\perp\mathbf 1$, fixed energy). That is a one-dimensional \emph{embedding} of the graph that respects adjacency. Sorting vertices by their Fiedler coordinate and cutting the order somewhere is the \emph{spectral ordering}; it is the workhorse heuristic for graph partitioning, sparse-matrix reordering (envelope/bandwidth reduction), and is the seed of Chapter~33.

\begin{whybox}
Why does minimizing $\sum_{(i,j)}(x_i-x_j)^2$ subject to $x\perp\mathbf1$ and $\norm{x}=1$ give a good cut? Because the objective is a sum over edges of squared coordinate gaps. To make it small you must give the two endpoints of almost every edge nearly equal coordinates -- so densely connected clumps collapse to nearly a single point on the line, and the only place coordinates are allowed to differ a lot is across the \emph{few} edges joining different clumps. The constraint $x\perp\mathbf1$ forbids the trivial all-equal solution and forces the clumps to spread to opposite signs. So the sign of $v_2$ is, heuristically, a label for ``which clump.''
\end{whybox}

\section{The Cheeger inequality}
\label{ch32:cheeger}

We now reach the central theorem. It makes the heuristic of the last box into a theorem with a constant, and it gives an algorithm.

\subsection{Conductance: the combinatorial quantity}

\begin{definition}[Cut, volume, conductance]
\label{ch32:conductance}
For $S\subseteq V$, write $S^c=V\setminus S$. The \emph{edge boundary} is $\partial S=\{(i,j)\in E : i\in S, j\in S^c\}$, with cut weight $w(\partial S)=\sum_{(i,j)\in\partial S}w_{ij}$. The \emph{volume} of $S$ is $\operatorname{vol}(S)=\sum_{i\in S}d_i$. The \emph{conductance} (a.k.a. Cheeger ratio, sparsity) of the cut is
\[
  \phi(S)=\frac{w(\partial S)}{\min\{\operatorname{vol}(S),\,\operatorname{vol}(S^c)\}},
\]
and the \emph{conductance of the graph} is $\phi(\Gph)=\min_{\emptyset\ne S\subsetneq V}\phi(S)$.
\end{definition}

Small $\phi(\Gph)$ means there is a way to split the vertices into two pieces, each of substantial volume, with only a few edges crossing -- a sparse cut, a bottleneck. Large $\phi(\Gph)$ means no such split exists; the graph is an \emph{expander} (Chapter~34). Computing $\phi(\Gph)$ exactly is NP-hard (it is essentially sparsest cut), so a polynomial-time two-sided estimate is valuable. Cheeger's inequality is that estimate.

\begin{intuitionbox}
Why divide by the smaller volume rather than by $\abs{S}$? Because we want to forbid trivial cuts that peel off a tiny set (a single vertex has cut $=d_i$ but is no real bottleneck). Normalizing by volume charges each side for the degree it contains, so a sparse cut must isolate a genuinely \emph{large} (high-volume) region cheaply. This is exactly the normalization the random walk sees, which is why $\mu_2$ -- not $\lambda_2$ -- is the eigenvalue that brackets $\phi$.
\end{intuitionbox}

\subsection{Statement}

\begin{theorem}[Cheeger inequality for graphs]
\label{ch32:cheeger-thm}
Let $\mu_2$ be the second-smallest eigenvalue of the normalized Laplacian $\mathcal{L}$ of a connected graph. Then
\[
  \frac{\mu_2}{2} \;\le\; \phi(\Gph) \;\le\; \sqrt{2\,\mu_2}.
\]
\end{theorem}

The two inequalities go in opposite directions and have completely different characters.
\begin{itemize}
  \item The \textbf{easy direction} $\tfrac{\mu_2}{2}\le\phi(\Gph)$ says ``a good cut forces a small eigenvalue.'' Its proof is a one-line Rayleigh-quotient test: plug the indicator of the sparse cut into the variational formula for $\mu_2$.
  \item The \textbf{hard direction} $\phi(\Gph)\le\sqrt{2\mu_2}$ says ``a small eigenvalue forces a good cut,'' and the proof is \emph{constructive}: it takes the eigenvector and produces a cut of conductance $\le\sqrt{2\mu_2}$. That construction is the spectral partitioning algorithm.
\end{itemize}

\begin{historybox}
The inequality is named for Jeff Cheeger, who proved the analytic version (a lower bound on the first nonzero eigenvalue of the Laplace--Beltrami operator on a Riemannian manifold in terms of an isoperimetric constant) in 1970. The discrete graph analogues were developed by Dodziuk, by Alon and Milman, and by Alon in the 1980s in the course of building and analyzing expander graphs. The constructive hard direction -- ``sweep the eigenvector'' -- is sometimes credited as the discrete Cheeger or the Cheeger rounding, and it is the theoretical heart of spectral clustering and of the Spielman--Teng program on nearly-linear-time Laplacian solvers (Chapter~41).
\end{historybox}

\subsection{Proof of the easy direction}

\begin{proof}[$\tfrac{\mu_2}{2}\le\phi(\Gph)$]
Let $S$ achieve $\phi(\Gph)$, with $\operatorname{vol}(S)\le\operatorname{vol}(S^c)$ (relabel if needed). We exhibit a test vector for the variational characterization of $\mu_2$. In the $y$-variables of Lemma~\ref{ch32:mu-le-2},
\[
  \mu_2 \;=\; \min_{\substack{y\ne 0\\ \sum_i d_i y_i = 0}}\ \frac{\sum_{(i,j)\in E} w_{ij}(y_i-y_j)^2}{\sum_i d_i\, y_i^2},
\]
where the constraint $\sum_i d_i y_i=0$ is ``$x\perp\mathbf{D}^{1/2}\mathbf1$,'' the orthogonality to the bottom eigenvector of $\mathcal{L}$. Choose
\[
  y_i \;=\;
  \begin{cases}
    \ \ \dfrac{1}{\operatorname{vol}(S)}, & i\in S,\\[2ex]
    -\dfrac{1}{\operatorname{vol}(S^c)}, & i\in S^c.
  \end{cases}
\]
This $y$ is constant on each side, so $(y_i-y_j)^2$ is nonzero only on boundary edges, each contributing $\big(\tfrac{1}{\operatorname{vol}(S)}+\tfrac{1}{\operatorname{vol}(S^c)}\big)^2$. The constraint holds: $\sum_i d_i y_i=\operatorname{vol}(S)\tfrac{1}{\operatorname{vol}(S)}-\operatorname{vol}(S^c)\tfrac{1}{\operatorname{vol}(S^c)}=0$. Compute both pieces:
\[
  \text{numerator}=w(\partial S)\Big(\tfrac{1}{\operatorname{vol}(S)}+\tfrac{1}{\operatorname{vol}(S^c)}\Big)^2,
  \qquad
  \text{denominator}=\tfrac{1}{\operatorname{vol}(S)}+\tfrac{1}{\operatorname{vol}(S^c)}
\]
(the denominator because $\sum_i d_i y_i^2=\operatorname{vol}(S)\tfrac{1}{\operatorname{vol}(S)^2}+\operatorname{vol}(S^c)\tfrac{1}{\operatorname{vol}(S^c)^2}$). Their ratio is
\[
  \mu_2 \le \frac{w(\partial S)\big(\tfrac{1}{\operatorname{vol}(S)}+\tfrac{1}{\operatorname{vol}(S^c)}\big)^2}
                  {\tfrac{1}{\operatorname{vol}(S)}+\tfrac{1}{\operatorname{vol}(S^c)}}
        = w(\partial S)\Big(\tfrac{1}{\operatorname{vol}(S)}+\tfrac{1}{\operatorname{vol}(S^c)}\Big)
        \le w(\partial S)\cdot\frac{2}{\operatorname{vol}(S)}
        = 2\,\phi(\Gph),
\]
using $\operatorname{vol}(S)\le\operatorname{vol}(S^c)$ so that $\tfrac1{\operatorname{vol}(S^c)}\le\tfrac1{\operatorname{vol}(S)}$. Rearranging, $\mu_2/2\le\phi(\Gph)$.
\end{proof}

\begin{intuitionbox}
The easy direction is just ``the cut indicator is a valid test signal.'' A sparse cut gives a near-piecewise-constant balanced signal of small roughness, and $\mu_2$ is by definition the minimum roughness, so $\mu_2$ is pushed down. The factor $2$ is the slack between a sharp $0/1$ indicator and the genuinely smoothest signal.
\end{intuitionbox}

\subsection{Proof sketch of the hard direction (the sweep)}

The hard direction turns the eigenvector into a cut. Let $f=v_2$ be the Fiedler vector of $\mathcal{L}$, written back in vertex space as $g=\mathbf{D}^{-1/2}f$ so that the relevant Rayleigh quotient is
\[
  \mu_2 \;=\; \frac{\sum_{(i,j)\in E} w_{ij}(g_i-g_j)^2}{\sum_i d_i g_i^2},
  \qquad \sum_i d_i g_i = 0 .
\]
We will produce a \emph{threshold cut}: sort the vertices by their value $g_i$, and consider the $n-1$ cuts $S_t=\{i : g_i \le \theta_t\}$ as the threshold $\theta_t$ sweeps through the sorted values. The claim is that the \emph{best} of these sweep cuts already has conductance $\le\sqrt{2\mu_2}$.

\begin{algorithm}[H]
\caption{\textsc{SpectralSweep} -- rounding a Fiedler vector to a cut}
\label{ch32:alg-sweep}
\KwIn{Connected graph $\Gph=(V,E,w)$, normalized Fiedler vector $f=v_2$ of $\mathcal{L}$.}
\KwOut{A vertex set $S^\star$ with $\phi(S^\star)\le\sqrt{2\,\mu_2}$.}
$g \leftarrow \mathbf{D}^{-1/2} f$\;
Sort vertices so that $g_{\sigma(1)}\le g_{\sigma(2)}\le\cdots\le g_{\sigma(n)}$\;
$\phi^\star \leftarrow +\infty$;\quad $S^\star\leftarrow\emptyset$\;
\For{$t=1$ \KwTo $n-1$}{
  $S_t \leftarrow \{\sigma(1),\dots,\sigma(t)\}$\tcp*{prefix in sorted order}
  Update $w(\partial S_t)$ and $\operatorname{vol}(S_t)$ incrementally from $S_{t-1}$\;
  $\phi(S_t) \leftarrow \dfrac{w(\partial S_t)}{\min\{\operatorname{vol}(S_t),\operatorname{vol}(V)-\operatorname{vol}(S_t)\}}$\;
  \If{$\phi(S_t)<\phi^\star$}{$\phi^\star\leftarrow\phi(S_t)$;\quad $S^\star\leftarrow S_t$\;}
}
\Return $S^\star$\;
\end{algorithm}

\begin{theorem}[Sweep guarantee, the hard direction]
\label{ch32:sweep-guarantee}
\textsc{SpectralSweep} returns a set $S^\star$ with $\phi(S^\star)\le\sqrt{2\,\mu_2}$. In particular $\phi(\Gph)\le\sqrt{2\,\mu_2}$.
\end{theorem}

\begin{proof}[Proof sketch]
The argument is a clever averaging. Shift and scale $g$ so that the median vertex (by volume) sits at $0$ and so that $\sum_i d_i g_i^2$ is normalized; one then splits $g$ into its positive part $g^+$ and negative part $g^-$ and works with whichever side carries at most half the volume -- call that vector $h$, supported on a set of volume $\le\tfrac12\operatorname{vol}(V)$, with $\dfrac{\sum_{(i,j)} w_{ij}(h_i-h_j)^2}{\sum_i d_i h_i^2}\le\mu_2$ still holding by a short calculation. Write $\rho$ for this Rayleigh quotient, so $\rho\le\mu_2$.

Now consider drawing a random threshold $\theta$ with density proportional to $2|\theta|$ over the range of $h$ (equivalently, threshold the values $h_i^2$ uniformly). For such a random threshold define the sweep set $S_\theta=\{i:h_i^2>\theta\}$ -- a prefix set in the sorted order. Two expectation computations finish it:
\begin{itemize}
  \item \emph{Expected boundary.} A boundary edge $(i,j)$ of $S_\theta$ is exactly one with $\theta$ between $h_i^2$ and $h_j^2$, which happens with probability $|h_i^2-h_j^2|=|h_i-h_j|\,|h_i+h_j|$. By Cauchy--Schwarz over edges,
  \[
    \mathbb{E}\,w(\partial S_\theta)
      =\sum_{(i,j)}w_{ij}|h_i-h_j||h_i+h_j|
      \le\sqrt{\textstyle\sum w_{ij}(h_i-h_j)^2}\ \sqrt{\textstyle\sum w_{ij}(h_i+h_j)^2}.
  \]
  The first factor is $\sqrt{\rho\sum d_i h_i^2}$. The second is $\le\sqrt{2\sum d_i h_i^2}$ by $(h_i+h_j)^2\le 2(h_i^2+h_j^2)$. So $\mathbb{E}\,w(\partial S_\theta)\le\sqrt{2\rho}\,\sum_i d_i h_i^2$.
  \item \emph{Expected volume.} A vertex $i$ lies in $S_\theta$ with probability $h_i^2/(\max_k h_k^2)$; scaling so $\max_k h_k^2=1$, $\mathbb{E}\operatorname{vol}(S_\theta)=\sum_i d_i h_i^2$, and every $S_\theta$ has volume $\le\tfrac12\operatorname{vol}(V)$ because $h$ is supported on the light side. Hence $\min\{\operatorname{vol}(S_\theta),\operatorname{vol}(S_\theta^c)\}=\operatorname{vol}(S_\theta)$.
\end{itemize}
Therefore
\[
  \mathbb{E}\,w(\partial S_\theta)
   \;\le\; \sqrt{2\rho}\;\mathbb{E}\operatorname{vol}(S_\theta).
\]
Because the left side is a nonnegative-weighted average of $w(\partial S_\theta)$ and the right side the same average of $\sqrt{2\rho}\operatorname{vol}(S_\theta)$, there must exist a particular threshold $\theta$ with $w(\partial S_\theta)\le\sqrt{2\rho}\operatorname{vol}(S_\theta)$, i.e. $\phi(S_\theta)\le\sqrt{2\rho}\le\sqrt{2\mu_2}$. That $S_\theta$ is one of the $n-1$ prefix sets the algorithm tries, so $S^\star$ is at least as good. (The averaging step is exactly the standard ``if $\mathbb{E}[A]\le\sqrt{2\rho}\,\mathbb{E}[B]$ with $B>0$ then some outcome has $A\le\sqrt{2\rho}\,B$'' applied to the two correlated quantities boundary and volume.)
\end{proof}

\begin{whybox}
The whole strength of the hard direction lives in the Cauchy--Schwarz step, which is where the square root $\sqrt{2\mu_2}$ -- rather than $\mu_2$ -- comes from. We could not avoid it: the gap between the two sides of Cheeger is real. The cycle $C_n$ realizes it. Its normalized eigenvalue is $\mu_2=1-\cos(2\pi/n)\approx \tfrac12(2\pi/n)^2=\Th(1/n^2)$, while its sparsest cut splits the cycle into two equal arcs: $2$ edges cross and each side has volume $n$, so $\phi(C_n)=2/n=\Th(1/n)$. Thus $\phi\approx\sqrt{\mu_2}$ up to constants, not $\phi\approx\mu_2$. The square-root loss is the price of getting a cut out of a real-valued eigenvector instead of an exact combinatorial optimum.
\end{whybox}

\begin{complexitybox}
Given the Fiedler vector, \textsc{SpectralSweep} costs $\Oh(m+n\log n)$: an $\Oh(n\log n)$ sort, then a single linear sweep maintaining $w(\partial S_t)$ and $\operatorname{vol}(S_t)$ incrementally (each vertex enters once; touch its incident edges, flipping each from ``crossing'' to ``internal'' or vice versa, in total $\Oh(m)$). Computing the Fiedler vector itself dominates: a few iterations of the Lanczos method or a Laplacian linear solve, $\widetilde{\Oh}(m)$ per matrix--vector product times the iteration count, or $\widetilde{\Oh}(m)$ overall with the nearly-linear-time solvers of Chapter~41. The guarantee $\phi(S^\star)\le\sqrt{2\mu_2}\le\sqrt{2}\sqrt{2\phi(\Gph)}=2\sqrt{\phi(\Gph)}$ makes spectral partitioning an $\Oh(1/\sqrt{\phi})$-approximation to sparsest cut -- weak in the worst case but excellent on graphs with a clear bottleneck.
\end{complexitybox}

\section{Worked example: Cheeger on a two-cluster graph}
\label{ch32:worked-cheeger}

Let us run the entire pipeline on a graph engineered to have one obvious bottleneck, so we can watch $\mu_2$, the Fiedler vector, and the sweep all point at the same cut.

Take six vertices in two triangles, $\{1,2,3\}$ and $\{4,5,6\}$, each a complete triangle, joined by a single bridge edge $3$--$4$. All weights $1$.

\begin{center}
\begin{tikzpicture}[every node/.style={circle,draw=nodeblue,thick,fill=blue!6,minimum size=8mm,inner sep=0pt},
  edge/.style={draw=edgegray,very thick}, bridge/.style={draw=trapred,very thick,dashed}]
  \node (v1) at (0,1.4) {$1$};
  \node (v2) at (0,-1.4) {$2$};
  \node (v3) at (1.6,0) {$3$};
  \node (v4) at (3.4,0) {$4$};
  \node (v5) at (5,1.4) {$5$};
  \node (v6) at (5,-1.4) {$6$};
  \draw[edge] (v1)--(v2) (v1)--(v3) (v2)--(v3);
  \draw[edge] (v4)--(v5) (v4)--(v6) (v5)--(v6);
  \draw[bridge] (v3)--(v4);
\end{tikzpicture}
\end{center}

Degrees: $d_1=d_2=2$, $d_3=3$, $d_4=3$, $d_5=d_6=2$, so $\operatorname{vol}(V)=14=2m$ with $m=7$. The natural cut $S=\{1,2,3\}$ has $w(\partial S)=1$ (the bridge) and $\operatorname{vol}(S)=7=\operatorname{vol}(S^c)$, hence
\[
  \phi(\{1,2,3\}) = \frac{1}{\min\{7,7\}} = \frac{1}{7}\approx 0.143,
\]
and one checks no other cut beats it, so $\phi(\Gph)=1/7$.

\textbf{The combinatorial Laplacian.}
\[
  \Lap=\begin{pmatrix}
   2 & -1 & -1 & 0 & 0 & 0\\
  -1 & 2 & -1 & 0 & 0 & 0\\
  -1 & -1 & 3 & -1 & 0 & 0\\
   0 & 0 & -1 & 3 & -1 & -1\\
   0 & 0 & 0 & -1 & 2 & -1\\
   0 & 0 & 0 & -1 & -1 & 2
  \end{pmatrix}.
\]
By the left--right symmetry (swap $1\leftrightarrow6$, $2\leftrightarrow5$, $3\leftrightarrow4$) every eigenvector is symmetric or antisymmetric. The Fiedler vector is the antisymmetric one that is constant on each triangle except for the bridge vertices: solving $\Lap v_2=\lambda_2 v_2$ with the ansatz $v_2\propto(b,b,a,-a,-b,-b)$, the equations decouple to a small $2\times2$ system in $a,b$, and we record the outcome:
\[
  \lambda_2 = \tfrac{5-\sqrt{17}}{2}\approx 0.438,\qquad v_2\approx(0.47,0.47,0.26,-0.26,-0.47,-0.47).
\]
The Fiedler vector is positive on the left triangle, negative on the right, and crosses zero exactly on the bridge. (The other eigenvalues are $\lambda_1=0$, then $\lambda_2\approx0.44$, and the rest clustered near $3$--$4.6$, the ``internal'' frequencies of the triangles.)

Let us actually carry out the $2\times2$ reduction, because it shows where that number comes from and why the symmetry ansatz is forced rather than guessed. Write the eigenvector as $v=(b,b,a,-a,-b,-b)$ -- the antisymmetric pattern that is constant on the two non-bridge vertices of each triangle. Two of the six eigen-equations are independent; the rest are copies by symmetry. Row~$1$ (vertex $1$, degree $2$, neighbors $2$ and $3$) reads
\[
  2b - b - a \;=\; \lambda b
  \qquad\Longrightarrow\qquad
  b - a \;=\; \lambda b .
\]
Row~$3$ (vertex $3$, degree $3$, neighbors $1,2,4$, with $v_4=-a$) reads
\[
  3a - b - b - (-a) \;=\; \lambda a
  \qquad\Longrightarrow\qquad
  4a - 2b \;=\; \lambda a .
\]
These are two linear equations in $(a,b)$; a nonzero solution needs the determinant of
\[
  \begin{pmatrix} 1-\lambda & -1 \\ -2 & 4-\lambda \end{pmatrix}
\]
to vanish, i.e. $(1-\lambda)(4-\lambda)-2=0$, which is $\lambda^2-5\lambda+2=0$. The two roots are $\lambda=\tfrac{5\pm\sqrt{17}}{2}$, giving $\lambda\approx 0.438$ and $\lambda\approx 4.56$. The smaller root, $\lambda_2=\tfrac{5-\sqrt{17}}{2}\approx 0.438$, is the antisymmetric Fiedler value; the larger is one of the internal triangle frequencies. Back-substituting $\lambda\approx0.438$ into $b-a=\lambda b$ gives $a=(1-\lambda)b\approx 0.562\,b$, so with the normalization $\norm v=1$ one recovers $v_2\approx(0.47,0.47,0.26,-0.26,-0.47,-0.47)$, matching the shape recorded above. The lesson is mechanical: a graph with a symmetry group lets you block-diagonalize the Laplacian by symmetry type and read each block's spectrum off a tiny matrix, no numerical solver required.

\textbf{Normalized version and Cheeger check.} With $\mathbf{D}^{-1/2}$ rescaling, the second normalized eigenvalue comes out to $\mu_2\approx 0.205$. Cheeger then brackets the true conductance $\phi(\Gph)=1/7\approx0.143$:
\[
  \underbrace{\tfrac{\mu_2}{2}\approx 0.102}_{\text{lower}}
  \ \le\
  \phi(\Gph)\approx 0.143
  \ \le\
  \underbrace{\sqrt{2\mu_2}\approx 0.640}_{\text{upper}}.
\]
Both inequalities hold, as promised, and the gap is exactly the square-root looseness of the hard direction.

\textbf{Running the sweep by hand.} Sort vertices by $g_i=\mathbf D^{-1/2}v_2$ value; ascending, the order groups the right triangle (negative coordinates) before the left (positive), so the meaningful break is between the two triangles. Sweep the prefix cuts:

\begin{center}
\begin{tabular}{clccc}
\toprule
$t$ & prefix $S_t$ (by sorted $g$) & $w(\partial S_t)$ & $\min\operatorname{vol}$ & $\phi(S_t)$\\
\midrule
1 & $\{5\}$              & $2$ & $2$ & $1.00$\\
2 & $\{5,6\}$            & $2$ & $4$ & $0.50$\\
3 & $\{4,5,6\}$          & $1$ & $7$ & $\mathbf{0.143}$\\
4 & $\{3,4,5,6\}$        & $2$ & $4$ & $0.50$\\
5 & $\{1,4,5,6\}$ \dots  & --  & --  & (worse)\\
\bottomrule
\end{tabular}
\end{center}

The minimum over the sweep is at $t=3$, the cut $\{4,5,6\}$ versus $\{1,2,3\}$, conductance $1/7$ -- the exact sparsest cut. The eigenvector found it. This is spectral partitioning working perfectly because the graph has a clean spectral gap (a big jump from $\lambda_2\approx0.44$ to the next eigenvalue), which is the regime where it shines. The full combinatorial spectrum makes the gap visible at a glance (Figure~\ref{ch32:spectrum-fig}).

\begin{figure}[H]
\centering
\begin{tikzpicture}[xscale=2.4]
  \draw[draw=edgegray,very thick,-{Stealth}] (-0.3,0) -- (5.0,0) node[right,font=\small]{$\lambda$};
  \foreach \x in {0,1,2,3,4}{\draw[draw=edgegray] (\x,0.06)--(\x,-0.06) node[below,font=\scriptsize]{$\x$};}
  % eigenvalues 0, 0.438, 3,3,3, 4.562
  \foreach \x/\lab in {0/{$\lambda_1$},0.438/{$\lambda_2$}}{
    \draw[draw=sourceblue,very thick] (\x,0)--(\x,0.45);
    \node[font=\scriptsize,sourceblue] at (\x,0.62){\lab};}
  \draw[draw=trapred,very thick] (3,0)--(3,0.45);
  \node[font=\scriptsize,trapred] at (3,0.62){$\lambda_{3,4,5}$};
  \draw[draw=trapred,very thick] (4.562,0)--(4.562,0.45);
  \node[font=\scriptsize,trapred] at (4.562,0.62){$\lambda_6$};
  \draw[<->,draw=quotientorange,thick] (0.438,-0.35)--(3,-0.35);
  \node[font=\scriptsize,quotientorange] at (1.7,-0.52){spectral gap};
\end{tikzpicture}
\caption{Combinatorial Laplacian spectrum of the two-triangle graph: $0,\ \tfrac{5-\sqrt{17}}{2}\!\approx\!0.44,\ 3,3,3,\ \tfrac{5+\sqrt{17}}{2}\!\approx\!4.56$. The wide gap between $\lambda_2$ and the cluster at $3$ signals that the graph splits cleanly into exactly two pieces.}
\label{ch32:spectrum-fig}
\end{figure}

It helps to see the embedding the Fiedler vector produces. Plot each vertex at its coordinate $g_i$ on a single axis: the six vertices collapse into two tight clumps with a wide gap straddling zero, and the sweep is nothing but ``scan left to right and try cutting in each gap.'' The widest gap -- the one between the two clumps -- is the cut the algorithm reports.

\begin{figure}[H]
\centering
\begin{tikzpicture}[>=Stealth,
  vnode_pos/.style={circle,draw=sourceblue,fill=sourceblue!15,thick,minimum size=6.5mm,inner sep=0pt,font=\scriptsize\bfseries},
  vnode_neg/.style={circle,draw=trapred,fill=trapred!15,thick,minimum size=6.5mm,inner sep=0pt,font=\scriptsize\bfseries}
]
  % Lane 1: Graph Topology and Nodal Partition
  \node[font=\footnotesize\bfseries\scshape,text=sourceblue!80!black,anchor=west] at (-5.5,2.6) {Lane 1: Combinatorial Graph $\Gph$ and Nodal Domains};
  \draw[gray!30,dashed] (-5.5,2.4) -- (5.5,2.4);

  % Left triangle (positive coordinates S)
  \node[vnode_pos] (u1) at (-3.2,1.8) {1};
  \node[vnode_pos] (u2) at (-3.2,0.2) {2};
  \node[vnode_pos] (u3) at (-1.6,1.0) {3};
  \draw[very thick,sourceblue] (u1)--(u2) (u2)--(u3) (u3)--(u1);

  % Right triangle (negative coordinates S^c)
  \node[vnode_neg] (u4) at (1.6,1.0) {4};
  \node[vnode_neg] (u5) at (3.2,1.8) {5};
  \node[vnode_neg] (u6) at (3.2,0.2) {6};
  \draw[very thick,trapred] (u4)--(u5) (u5)--(u6) (u6)--(u4);

  % Bridge edge
  \draw[line width=1.5pt,edgegray] (u3) -- node[midway,above=2pt,font=\tiny\bfseries,black] {bridge} (u4);

  % Cut line separating S and S^c
  \draw[dashed,line width=1.5pt,quotientorange] (0,2.2) -- (0,-0.2) node[below=2pt,font=\scriptsize\bfseries,quotientorange] {Cheeger cut $\partial S = \{3,4\}$};

  \node[font=\scriptsize,sourceblue,anchor=east] at (-3.6,1.0) {Cluster $S$ ($v_2 > 0$)};
  \node[font=\scriptsize,trapred,anchor=west] at (3.6,1.0) {Cluster $S^c$ ($v_2 < 0$)};

  % Divider
  \draw[gray!40,thick] (-5.5,-0.6) -- (5.5,-0.6);

  % Lane 2: 1D Spectral Embedding and Sweep
  \node[font=\footnotesize\bfseries\scshape,text=quotientorange!80!black,anchor=west] at (-5.5,-1.0) {Lane 2: 1D Spectral Embedding $g_i = \mathbf{D}^{-1/2}v_2(i)$ and Sweep Cut};
  \draw[gray!30,dashed] (-5.5,-1.2) -- (5.5,-1.2);

  % Coordinate axis
  \draw[thick,edgegray,->] (-5.0,-2.6) -- (5.0,-2.6) node[right=2pt,font=\small] {$g_i$};
  \foreach \x/\val in {-4/-0.33, -2/-0.15, 0/0, 2/+0.15, 4/+0.33} {
    \draw[thick,edgegray] (\x,-2.5) -- (\x,-2.7) node[below=2pt,font=\tiny] {$\val$};
  }

  % Plotted vertices
  \node[vnode_neg] at (-4.0,-2.0) {5};
  \node[vnode_neg] at (-4.0,-3.3) {6};
  \node[vnode_neg] at (-2.0,-2.0) {4};

  \node[vnode_pos] at (2.0,-2.0) {3};
  \node[vnode_pos] at (4.0,-2.0) {1};
  \node[vnode_pos] at (4.0,-3.3) {2};

  % Cut threshold at 0
  \draw[dashed,line width=1.5pt,quotientorange] (0,-1.5) -- (0,-3.7) node[below=2pt,font=\scriptsize\bfseries,quotientorange] {Optimal threshold $\theta^* = 0$};

  \node[font=\tiny,align=center,fill=white,inner sep=2pt] at (0,-1.6) {Minimal conductance: $\phi(S) = \frac{1}{7} \approx 0.143$};
\end{tikzpicture}
\caption{The Cheeger nodal cut along the Fiedler vector for the two-triangle graph. \textbf{Top (Lane~1):} The graph vertices partition into two nodal domains according to the sign of the Fiedler vector coordinates: $S=\{1,2,3\}$ (positive, blue) and $S^c=\{4,5,6\}$ (negative, red), severing only the bridge edge $\{3,4\}$. \textbf{Bottom (Lane~2):} Embedding vertices on the line at $g_i=\mathbf{D}^{-1/2}v_2(i)$ clumps the two clusters on opposite sides of zero. Sweeping the threshold across the gap isolates the exact sparsest cut with conductance $\phi = 1/7$.}
\label{ch32:fig-cheeger-nodal}
\end{figure}

The dashed line is the winning threshold: every edge except the bridge has both endpoints on the same side, so only one edge crosses, and that is the conductance-$1/7$ cut the table found at $t=3$. This picture is the whole idea of spectral partitioning in one image -- the eigenvector turns a combinatorial cut into a one-dimensional gap-finding problem.

\begin{labbox}
Redo this with the bridge \emph{doubled} (weight $2$ on edge $3$--$4$, or a second parallel bridge). Recompute $\partial S$ for $S=\{1,2,3\}$, the new $\operatorname{vol}$, and predict the direction $\lambda_2$ moves before computing it. (It must increase, by edge monotonicity, Proposition~\ref{ch32:monotone}.) Then verify the sweep still picks $\{1,2,3\}$ but with larger conductance. This is the cleanest hands-on way to feel ``more bridge $\Rightarrow$ less of a bottleneck $\Rightarrow$ larger $\lambda_2$.''
\end{labbox}

\begin{applicationbox}
Spectral partitioning via the sweep is the standard first cut in: \emph{image segmentation} (the normalized cuts algorithm of Shi and Malik is exactly $\mathcal{L}$'s Fiedler vector on a pixel-similarity graph); \emph{scientific computing} (partitioning a finite-element mesh across processors to minimize inter-processor communication, where the mesh graph's bottleneck is the seam to cut); \emph{community detection} in social and biological networks; and \emph{sparse-matrix reordering} to reduce fill-in. In each case the win is the same: a hard combinatorial cut problem is replaced by an eigenvector computation plus a linear sweep, with a Cheeger guarantee on the result.
\end{applicationbox}

\section{Higher eigenvalues and many-way partitions}
\label{ch32:higher}

The story does not stop at $\lambda_2$. The eigenvectors $v_2,v_3,\dots,v_k$ together give a $(k-1)$-dimensional embedding of the graph -- map vertex $i$ to the point $(v_2(i),v_3(i),\dots,v_k(i))$ -- and clustering that point cloud (e.g. with $k$-means) is the $k$-way spectral clustering of Chapter~33. The theoretical justification is a higher-order Cheeger inequality.

\begin{theorem}[Higher-order Cheeger, informal]
\label{ch32:higher-cheeger}
Let $\rho_k(\Gph)=\min\big\{\max_{1\le i\le k}\phi(S_i)\big\}$ over partitions of $V$ into $k$ nonempty disjoint sets $S_1,\dots,S_k$ (the $k$-way expansion). Then there are universal constants such that
\[
  \frac{\mu_k}{2}\ \le\ \rho_k(\Gph)\ \le\ C\,k^2\,\sqrt{\mu_k},
\]
and a refined bound replaces $k^2$ by $\Oh(k)$ using $\Oh(k)$ eigenvectors. Consequently $\Gph$ has a meaningful $k$-clustering exactly when there is a large gap $\mu_{k+1}\gg\mu_k$ in the spectrum.
\end{theorem}

We state this without proof (it is due to Lee, Oveis Gharan, and Trevisan, 2012); the takeaway is the \emph{spectral-gap heuristic}: \emph{the number of clusters to ask for is the index $k$ at which the eigenvalue sequence $\mu_2,\mu_3,\dots$ first jumps}. In the two-triangle example there is a big jump after $\mu_2$, signalling exactly two clusters. We return to this in the next chapter.

\begin{connectionbox}
The eigenvalue sequence $0=\mu_1\le\mu_2\le\cdots$ is the same object that governs the \emph{mixing time} of the random walk (Chapter~36) -- the walk mixes in roughly $\Oh(\tfrac{1}{\mu_2}\log n)$ steps, so a small $\mu_2$ (a bottleneck) means slow mixing, which is the probabilistic shadow of Cheeger. It is also the quantity expander constructions (Chapter~34) are built to control: an expander is precisely a sparse graph with $\mu_2$ bounded away from $0$ (equivalently $\phi(\Gph)=\Omega(1)$), and Cheeger is what certifies the two notions agree. The thread runs straight on to the Laplacian solvers of Chapter~41, whose convergence is governed by the spectrum.
\end{connectionbox}

\section{Stability: how much does the spectrum move?}
\label{ch32:stability}

A reason the spectrum is useful in practice is that it is \emph{stable}: small changes to the graph make small changes to the eigenvalues. This is Weyl's inequality applied to the Laplacian.

\begin{proposition}[Weyl stability]
\label{ch32:weyl}
If $\Lap$ and $\Lap'$ are Laplacians of two weighted graphs on the same vertex set, then for every $k$,
\[
  \big|\lambda_k(\Lap)-\lambda_k(\Lap')\big|\ \le\ \norm{\Lap-\Lap'}_2,
\]
the spectral norm of the difference. In particular, adding or deleting a single unit edge changes each $\lambda_k$ by at most $2$ (the spectral norm of one edge Laplacian $\Lap_e$ is $2$).
\end{proposition}

\begin{proof}
Weyl's inequality for symmetric matrices states $\lambda_k(M+E)\le\lambda_k(M)+\lambda_{\max}(E)$ and $\lambda_k(M+E)\ge\lambda_k(M)+\lambda_{\min}(E)$; combining with $E=\Lap'-\Lap$ gives $|\lambda_k(\Lap')-\lambda_k(\Lap)|\le\max(|\lambda_{\max}(E)|,|\lambda_{\min}(E)|)=\norm{E}_2$. For a single edge $e=(p,q)$, $\Lap_e=(\mathbf e_p-\mathbf e_q)(\mathbf e_p-\mathbf e_q)^\top$ has one nonzero eigenvalue $\norm{\mathbf e_p-\mathbf e_q}^2=2$.
\end{proof}

This robustness is what makes spectral methods workable on noisy real-world graphs and what underlies \emph{spectral sparsification} (Chapter~41): one can replace a dense graph by a sparse weighted one whose Laplacian is spectrally close, $(1-\eps)\Lap\preceq\widetilde\Lap\preceq(1+\eps)\Lap$, and every eigenvalue -- hence every Cheeger estimate, every cut value up to $1\pm\eps$, every effective resistance -- is preserved. The spectrum is the right currency precisely because it changes continuously with the graph.

\begin{trapbox}
Stability of \emph{eigenvalues} does not imply stability of \emph{eigenvectors}. When $\mu_2$ and $\mu_3$ are nearly equal (a small spectral gap), the Fiedler vector can swing wildly under a tiny perturbation, because the relevant eigenspace is nearly two-dimensional and the choice of axis within it is ill-conditioned. The Davis--Kahan $\sin\Theta$ theorem makes this precise: eigenvector stability is governed by the \emph{gap} to neighboring eigenvalues, not by the perturbation size alone. So always check the spectral gap before trusting a single Fiedler cut; if the gap is small, the graph does not have a single well-defined bottleneck and you should ask for more clusters.
\end{trapbox}

\section{Summary}
\label{ch32:summary}

The dictionary we built:
\begin{itemize}
  \item $\lambda_1=\mu_1=0$ always (the constant vector); multiplicity of $0$ counts connected components (Theorem~\ref{ch32:components}).
  \item $0\le\mu_k\le 2$; $\mu_n=2$ iff a component is bipartite (Lemma~\ref{ch32:mu-le-2}).
  \item $\lambda_2$, the algebraic connectivity / Fiedler value, measures how robustly the graph hangs together: monotone in edges, $\le$ vertex connectivity, $=0$ iff disconnected (Definition~\ref{ch32:def-fiedler}, Proposition~\ref{ch32:monotone}).
  \item Cheeger: $\tfrac{\mu_2}{2}\le\phi(\Gph)\le\sqrt{2\mu_2}$ (Theorem~\ref{ch32:cheeger-thm}); the upper bound is constructive via the eigenvector sweep (Theorem~\ref{ch32:sweep-guarantee}, Algorithm~\ref{ch32:alg-sweep}), giving a polynomial-time sparse-cut approximation.
  \item Higher eigenvalues give $k$-way partitions; the spectral gap tells you how many clusters there are (Theorem~\ref{ch32:higher-cheeger}).
  \item The spectrum is Weyl-stable under edge perturbations and sparsification, but eigenvectors are only stable when the spectral gap is large (Proposition~\ref{ch32:weyl}).
\end{itemize}

The two-triangle worked example tied it together: $\mu_2$ small, a clean spectral gap, a Fiedler vector that changes sign exactly across the bridge, and a sweep that recovers the exact sparsest cut. The next chapter turns this from one cut into a clustering pipeline; Chapter~34 studies the graphs where $\phi$ is bounded below (expanders); and Chapter~41 makes the eigenvector and solver computations nearly linear time.

\begin{connectionbox}
\textbf{What This Unlocked:}
Spectral graph theory turned NP-hard graph partitioning into continuous eigenvalue optimization.
The discrete Cheeger inequality guaranteed that sorting coordinates of the Fiedler vector
yields an $\Oh(\sqrt{\mu_2})$ approximation to the sparsest cut in polynomial time.
This unlocked spectral graph bisection, image segmentation (Normalized Cuts by Shi and Malik 2000),
random-walk mixing time analysis, and the foundational theory of expander graphs.

\textbf{Honest Downsides and Limitations:}
\begin{itemize}
  \item \emph{The square-root gap:} Cheeger's inequality has an inherent quadratic loss:
        $\tfrac{\mu_2}{2} \le \phi(\Gph) \le \sqrt{2\mu_2}$. When the true conductance $\phi$ is
        small (e.g. on cycles $C_n$ or hypercubes $Q_d$), the spectral sweep guarantees only
        an $\Oh(\sqrt{\phi})$ cut rather than an $\Oh(\phi)$ cut. The Arora--Rao--Vazirani (2009)
        algorithm uses semidefinite programming and $\ell_2^2$ embeddings to bridge this gap to $\Oh(\sqrt{\log n})$.
  \item \emph{Multi-cluster blind spot:} A single Fiedler vector partitions a graph into only
        two pieces. It cannot natively separate graphs with $k > 2$ clusters or complex hierarchical
        communities without recursive sweeps or multiple eigenvectors.
\end{itemize}

\textbf{Where We Go Next:}
How do we partition a graph into $k \ge 3$ natural clusters simultaneously?
In \textbf{Chapter~33 (Spectral Clustering)}, we embed vertices into $\R^k$ using the top $k$
Laplacian eigenvectors and cluster them via $k$-means.
And what happens when $\mu_2 \ge c > 0$ is bounded away from zero regardless of graph size?
Such graphs have no bottlenecks and mix instantly: they are \textbf{Chapter~34's Expander Graphs},
the bedrock of derandomization, coding theory, and modern theoretical computer science.
\end{connectionbox}

\section{Exercises}
\label{ch32:exercises}

\begin{exercise}[Star and its spectrum]
Compute the Laplacian spectrum of the star $K_{1,n-1}$ (one center joined to $n-1$ leaves). Show $\lambda_2=1$ and $\lambda_n=n$, with $\lambda=1$ of multiplicity $n-2$. Interpret $\lambda_2=1$: how sparse is the sparsest cut, and does Cheeger's bracket contain it?
\end{exercise}

\begin{exercise}[Disjoint union]
Prove that if $\Gph$ is the disjoint union of $\Gph_1$ and $\Gph_2$, then the Laplacian spectrum of $\Gph$ is the multiset union of the two spectra. Deduce the component-counting theorem (Theorem~\ref{ch32:components}) as a corollary.
\end{exercise}

\begin{exercise}[Edge monotonicity is tight]
Give a graph and an edge whose addition raises $\lambda_2$ by exactly the maximum allowed by Weyl ($+2$), and another whose addition raises $\lambda_2$ by almost nothing. Explain combinatorially what distinguishes the two edges.
\end{exercise}

\begin{exercise}[Bipartite double cover]
For any graph $\Gph$, its bipartite double cover $\Gph\times K_2$ has vertex set $V\times\{0,1\}$ with edges $((u,0),(v,1))$ for each $uv\in E$. Show its normalized spectrum is $\{\mu_k\}\cup\{2-\mu_k\}$. Conclude $\mu_n(\Gph\times K_2)=2$ always, consistent with Lemma~\ref{ch32:mu-le-2}.
\end{exercise}

\begin{exercise}[Cheeger lower bound by hand]
Take the path $P_4$ (Section~\ref{ch32:p4}). Compute $\phi(P_4)$ exactly (try all cuts), compute $\mu_2$, and verify $\tfrac{\mu_2}{2}\le\phi\le\sqrt{2\mu_2}$ numerically. Which side of Cheeger is closer to tight here, and why?
\end{exercise}

\begin{exercise}[Implement the sweep]
Implement Algorithm~\ref{ch32:alg-sweep}. Feed it a graph that is two cliques $K_t$ joined by a single edge. Verify that the returned cut is the bridge cut and that its conductance matches $1/(t(t-1)+1)$ up to your normalization. Then make the bridge $r$ parallel edges and watch $\phi$ and $\mu_2$ both rise.
\end{exercise}

\begin{exercise}[Break the sweep]
Construct a graph where the \emph{global} sparsest cut is \emph{not} a prefix of the Fiedler ordering, so the single sweep misses it. (Hint: nest two bottlenecks at different scales, or use a graph whose Fiedler vector has its small spectral gap so $v_2$ mixes two cuts.) Estimate how far off the sweep's conductance is, and relate it to the spectral gap $\mu_3-\mu_2$.
\end{exercise}

\begin{exercise}[Fiedler nodal domains]
Prove that on a \emph{tree}, the Fiedler vector $v_2$ has the ``one sign change'' property: the vertices with $v_2(i)>0$ form a connected subtree and likewise for $v_2(i)<0$. (This is Fiedler's nodal-domain theorem for trees; use the Rayleigh-quotient characterization and an exchange argument.) Show by example it can fail on a general graph.
\end{exercise}

\begin{exercise}[Algebraic vs.\ vertex connectivity]
Verify Fiedler's bound $\lambda_2\le\kappa(\Gph)$ on three graphs: $C_n$, $K_{3,3}$, and the Petersen graph. For which is the bound tight? Then exhibit a graph where $\lambda_2$ is much smaller than $\kappa$, demonstrating that algebraic connectivity can badly underestimate vertex connectivity.
\end{exercise}

\begin{exercise}[Weyl in action]
Let $\Gph$ be the cycle $C_{100}$. Predict, using Proposition~\ref{ch32:weyl}, the maximum possible change in $\lambda_2$ when you add a single ``shortcut'' chord. Then compute the actual change for the chord joining antipodal vertices and explain why it is so much larger than for a short chord -- connect this to the chord's effect on the bottleneck.
\end{exercise}

\begin{exercise}[Normalized vs.\ combinatorial]
Build a graph with a highly skewed degree distribution (one high-degree hub, many low-degree leaves) where the combinatorial Fiedler vector and the normalized Fiedler vector suggest \emph{different} cuts. Explain which one respects volume and why the normalized version is the right one for conductance.
\end{exercise}

\begin{exercise}[Cheeger constant of the hypercube]
The hypercube $Q_d$ has $2^d$ vertices and is $d$-regular with $\mu_2=2/d$. Use Cheeger to bound $\phi(Q_d)$ from both sides, then compare with the exact isoperimetric fact $\phi(Q_d)=1/d$ (the cut splitting on one coordinate). Which Cheeger direction is tight, and what does that say about the square-root loss for this graph?
\end{exercise}

\begin{exercise}[Two-triangle spectrum by symmetry]
Redo the $2\times2$ reduction of Section~\ref{ch32:worked-cheeger} for the \emph{symmetric} eigenvectors $v\propto(b,b,a,a,b,b)$ (constant across the bridge instead of flipping). Write down the two eigen-equations, find the resulting quadratic, and confirm that one of its roots is $\lambda_1=0$ (the all-ones direction) so the symmetric block contributes the trivial eigenvalue plus one nontrivial value. Combined with the antisymmetric block you computed in the text, you now have four of the six eigenvalues from two $2\times2$ determinants; identify where the remaining two (the pure within-triangle modes orthogonal to both $a$ and $b$ patterns) come from.
\end{exercise}

\begin{exercise}[Conductance versus the unnormalized cut ratio]
Define the unnormalized ratio $h(S)=w(\partial S)/\min\{\abs S,\abs{S^c}\}$ (counting vertices, not volume). Show $h$ and $\phi$ agree up to factors of the maximum and minimum degree, and exhibit a graph with a hub where $h$ and $\phi$ are minimized by \emph{different} sets. Explain, using the easy-direction proof, why $\phi$ -- not $h$ -- is the one Cheeger's inequality is stated for.
\end{exercise}
